% Generated draft: requires mathematical and source review.
\documentclass[11pt]{article}
\usepackage[margin=25mm]{geometry}
\usepackage{fontspec}
\setmainfont{Latin Modern Roman}
\usepackage{amsmath,amssymb,mathtools,unicode-math}
\setmathfont{Latin Modern Math}
\usepackage{enumitem,needspace,xurl,hyperref}
\setlength{\parindent}{0pt}
\setlength{\parskip}{5pt}
\setlength{\emergencystretch}{3em}
\newcommand{\N}{\mathbb{N}}
\newcommand{\Z}{\mathbb{Z}}
\newcommand{\Q}{\mathbb{Q}}
\newcommand{\R}{\mathbb{R}}
\newcommand{\C}{\mathbb{C}}
\newcommand{\F}{\mathbb{F}}
\newcommand{\A}{\mathbb{A}}
\newcommand{\PP}{\mathbb{P}}
\newcommand{\E}{\mathbb{E}}
\newcommand{\eps}{\varepsilon}
\newcommand{\dd}{\,\mathrm{d}}
\begin{document}
\section*{UnsolvedMath ID: 30005805}
\textbf{Source problem code:} catalogue entry
\subsection{Definitions and mathematical statement}
For each integer \(n\geq 4\), let \(V_n=\{v_1,\ldots,v_n\}\). Define \(\mathcal T_n\) to be the family of labeled stacked triangulations constructed from the complete graph \(K_4\) on \(v_1,\ldots,v_4\) as follows. For each \(i=4,\ldots,n-1\), choose one triangular face of the current triangulation, insert \(v_{i+1}\) into that face, and join it to the three vertices bounding that face. Faces are understood combinatorially on the sphere: there are \(2i-4\) faces at this stage. In a plane representation this count includes the unbounded face. Members of \(\mathcal T_n\) are the resulting labeled graphs.

For a graph \(G\) on \(V_n\) and a set \(P\subset\R^2\) of \(n\) distinct points, a crossing-free straight-line embedding of \(G\) on \(P\) is a bijection \(\phi:V_n\to P\) such that each edge \(uv\) is represented by the segment \([\phi(u),\phi(v)]\), no vertex image lies in the relative interior of an edge segment, and distinct edge segments intersect only at the image of a common endpoint, if they have one.

Let \(G_1,G_2\) be independent random graphs, each uniformly distributed on \(\mathcal T_n\). Define
\[
 p_n=\Pr\!\left(\begin{array}{l}
 \text{there exists }P\subset\R^2,\ |P|=n,\text{ such that both }G_1\text{ and }G_2\\
 \text{have crossing-free straight-line embeddings on }P
 \end{array}\right).
\]
The two embeddings may use different bijections \(\phi_1,\phi_2:V_n\to P\); equally labeled vertices are not required to occupy the same point. Crossing-freeness is required separately for each drawing, with no restriction on intersections between edges from different drawings. No general-position restriction on \(P\) is imposed in the supplied reconstructed statement.

Determine the asymptotic behavior of \(p_n\) as \(n\to\infty\), including its limit if one exists. The source asks for the probability ``for large \(n\)'' without specifying a desired asymptotic precision. The distribution and common-point-set formulation above follow the supplied reconstructed statement; the supplied original excerpt refers instead to the distribution in the proofs of Theorems 1 and 2 without reproducing its definition. Those conventions require source review. The registry labels this problem open and the reconstruction \texttt{corrected\_verified}; these are supplied curation labels, not independent verification.
\subsection{Short English statement}
Independently choose two uniformly random labeled stacked triangulations on \(n\) vertices from the supplied construction. How does the probability that both can be drawn with straight edges and no crossings on one common set of \(n\) plane points behave as \(n\to\infty\)? Each graph may assign its vertices to those points differently.
\subsection{Sources}
\begin{itemize}
\item Supplied local UnsolvedMath registry record, permanent ID \(30005805\), title \emph{Simultaneous Embeddings of Random Planar Graphs}, exact source problem code \texttt{catalogue entry}. The record attributes the problem to the Oberwolfach Reports Open Problems collection and gives proposed year 2024; no individual proposer is supplied.
\item \emph{Discrete Geometry} (2024), Oberwolfach Reports, pp.\ 137--202, DOI \texttt{10.4171/owr/2024/3}. Supplied report reference and literature source: \url{https://doi.org/10.4171/owr/2024/3}. The retained original passage asks about simultaneous embeddings for two independent graphs from the distribution used in the proofs of Theorems 1 and 2, and mentions its relevance to Problem 1 without supplying that problem.
\item \emph{A Logarithmic Bound for Simultaneous Embeddings of Planar Graphs}. Supplied literature source: \url{https://doi.org/10.1007/s00454-024-00665-7}.
\item Literature assessment dated exactly \texttt{2026-08-22}: supplied dated curation, not independently verified. It labels the problem open, identifies the two-random-graph question as Problem 6.1 of the 2024 primary paper, and reports a non-embeddability result for about \(3\log_2(n)\) independent stacked triangulations, rather than the requested two-graph limit. The record also states that its reconstruction was checked against the official Oberwolfach report; that verification claim is supplied metadata and has not been independently checked here.
\end{itemize}
\end{document}
